% titel: Blatt 8, A3: Enthalpieänderung beim Erwärmen (Eis $\to$ Dampf)
% skript: Mortimer Kap. 11 (Thermochemie, kalorimetrische Zustandsänderungen)

% Teil I: Werkzeuge, kein Lösungsweg.

Um einen Stoff von einer Temperatur durch einen Phasenwechsel in eine andere
zu bringen, addiert man die Beiträge der einzelnen Strecken. Es gibt zwei
unterschiedliche Arten von Beiträgen.

\begin{werkzeugkasten}[Erwärmen und Phasenwechsel]
Innerhalb eines Phasenbereichs gilt (keine Phasenänderung):
\[
  q = m \cdot c \cdot \Delta T .
\]
Dabei sind $m$ die Masse, $c$ die spezifische (oft massenbezogene)
Wärmekapazität und $\Delta T$ die Temperaturdifferenz.

An einer Phasenwechsel-Temperatur $T$ ist die Temperatur konstant, man
führt Schmelz- bzw. Verdampfungsenthalpie zu:
\[
  q = n \cdot \Delta H_{\text{fus}}\quad\text{bzw.}\quad
  q = n \cdot \Delta H_{\text{vap}} .
\]
\textit{Skript:} Mortimer Kap. 11; die spezifischen Wärmekapazitäten sind
je Phase verschieden (fester, flüssiger, gasförmiger Zustand).
\end{werkzeugkasten}

\begin{achtung}[Darauf achten]
\begin{itemize}
  \item Die $c$-Werte gelten jeweils \emph{einmal pro Phase}; zwischen den
        Umwandlungspunkten darf die Temperatur nicht über die Sprungstelle
        gezogen werden (erst bis 0\,°C, dann schmelzen, dann bis 100\,°C, ...).
  \item Einheiten: $J\,g^{-1}K^{-1}$ und $kJ\,mol^{-1}$ mischen sich --- erst
        einheitlich umrechnen (mol $\rightarrow$ g, kJ $\rightarrow$ J oder umgekehrt).
  \item Es gibt \emph{drei} Erwärmungsstufen und \emph{zwei} Phasenwechsel.
  \item Teilchenanzahl: 1.00 mol H$_2$O entspricht einer konkreten Masse.
\end{itemize}
\end{achtung}
